\clearpage\thispagestyle{empty}
\pdfbookmark[0]{Quellen zu Gauß an von Lindenau, 30. August 1808}{lindenau1808aug-sources}
\begin{center}\textbf{Quellen zu Gauß an von Lindenau, 30. August 1808}\end{center}
\noindent Diese KI-Übertragung in heutiges Deutsch folgt Carl Friedrich Gauß’ deutschem Text in den \textit{Werken}, Band~VI (Göttingen, 1874), S.~146--147, einschließlich der Anmerkung zum Sextanten. Der Abschnitt steht auf S.~271--273 des Beitrags XXXIV, \textit{Auszug aus einem Schreiben des Herrn Professors Gauss}, im \href{https://www.digitale-sammlungen.de/en/view/bsb10538609?page=293}{Septemberheft 1808} der von Franz Xaver von Zach herausgegebenen \textit{Monatlichen Correspondenz zur Beförderung der Erd- und Himmels-Kunde}, Band~18, S.~269--273. Die Anmerkung der Zeitschrift auf S.~271 kündigt die Abhandlung zur sphärischen Astronomie im \href{https://www.digitale-sammlungen.de/en/view/bsb10538609?page=299}{folgenden Heft} an. Die Werke drucken \textit{Reispiel}, die Zeitschrift \textit{Beyspiel}. \href{https://gdz.sub.uni-goettingen.de/id/PPN236007467}{Historischer deutscher Text in den Werken}.
\par\smallskip\noindent Übertragung und Satz wurden mit KI erstellt.
\clearpage
